\documentclass[25pt, a0paper, portrait, margin=0mm, innermargin=15mm,
               blockverticalspace=15mm, colspace=15mm, subcolspace=8mm]{tikzposter}

\usepackage[utf8]{inputenc}
\usepackage[spanish,es-nodecimaldot]{babel}
\usepackage[T1]{fontenc}
\usepackage{lmodern}
\usepackage{amsmath, amssymb}
\usepackage{graphicx}
\usepackage{booktabs}
\usepackage{array}
\usepackage{tikz}
\usetikzlibrary{shapes.geometric, arrows.meta, positioning, calc, backgrounds, fit}
\usepackage{xcolor}
\usepackage{qrcode}
\usepackage{multicol}
\usepackage{siunitx}

% =====================================================================
% Paleta de colores estilo astronomia
% =====================================================================
\definecolor{astrodark}{HTML}{0A0E27}
\definecolor{astroblue}{HTML}{1B2A5C}
\definecolor{astrogold}{HTML}{FFD700}
\definecolor{astrocyan}{HTML}{4FC3F7}
\definecolor{astrolight}{HTML}{E8EEF7}
\definecolor{clumpyred}{HTML}{E63946}
\definecolor{noclumpygreen}{HTML}{2A9D8F}

% =====================================================================
% Estilo personalizado del poster
% =====================================================================
\definetitlestyle{astrotitle}{
    width=1150mm, roundedcorners=15, linewidth=3pt,
    innersep=15pt, titletotopverticalspace=15mm,
    titletoblockverticalspace=15mm
}{
    \begin{scope}[line width=\titlelinewidth,
                  rounded corners=\titleroundedcorners]
        \fill[color=astrodark] (\titleposleft, \titleposbottom)
            rectangle (\titleposright, \titlepostop);
        \draw[color=astrogold] (\titleposleft, \titleposbottom)
            rectangle (\titleposright, \titlepostop);
    \end{scope}
}
\usetitlestyle{astrotitle}

\defineblockstyle{astroblock}{
    titlewidthscale=1, bodywidthscale=1, titleleft,
    titleoffsetx=0pt, titleoffsety=0pt,
    bodyoffsetx=0pt, bodyoffsety=0pt,
    bodyverticalshift=0pt, roundedcorners=8, linewidth=2pt,
    titleinnersep=8mm, bodyinnersep=8mm
}{
    \draw[color=astroblue, fill=astrolight, line width=\blocklinewidth,
          rounded corners=\blockroundedcorners] (blockbody.south west)
        rectangle (blockbody.north east);
    \ifBlockHasTitle
        \draw[color=astroblue, fill=astroblue, line width=\blocklinewidth,
              rounded corners=\blockroundedcorners] (blocktitle.south west)
            rectangle (blocktitle.north east);
    \fi
}
\useblockstyle{astroblock}

\colorlet{blocktitlebgcolor}{astroblue}
\colorlet{blocktitlefgcolor}{astrogold}
\colorlet{blockbodybgcolor}{astrolight}
\colorlet{blockbodyfgcolor}{astrodark}
\colorlet{backgroundcolor}{white}

% =====================================================================
% Metadata del poster
% =====================================================================
\title{\parbox{\linewidth}{\centering
    \color{astrogold}\fontsize{80}{95}\selectfont\bfseries
    Clasificacion Automatica de Galaxias Clumpy \\[8pt]
    \color{astrocyan}\fontsize{55}{65}\selectfont\bfseries
    en Imagenes JWST usando Deep Learning Hibrido
}}
\author{\color{white}\Large\bfseries
    Autor: [TU NOMBRE] \quad$|$\quad
    Guia: [NOMBRE DEL GUIA] \quad$|$\quad
    [DEPARTAMENTO / CARRERA]
}
\institute{\color{astrocyan}\normalsize
    [UNIVERSIDAD] \quad$\cdot$\quad Proyecto Semestral de Deep Learning \quad$\cdot$\quad 2026
}
\titlegraphic{
    % Espacio para logo universidad. Reemplazar path por logo real.
    % \includegraphics[width=8cm]{logo_universidad.png}
    \color{astrogold}\Huge$\bigstar$
}

\begin{document}
\maketitle[titletotopverticalspace=15mm]

% =====================================================================
% BLOQUE SUPERIOR: RESUMEN + METRICAS DESTACADAS
% =====================================================================
\begin{columns}
\column{0.65}
\block{Resumen}{
\large
Presentamos un clasificador binario para detectar galaxias \textbf{clumpy}
(con morfologia grumosa, asociada a formacion estelar intensa) en imagenes
del telescopio \textbf{James Webb (JWST)}. La version final (\textbf{v2.1})
combina un backbone convolucional moderno (\textbf{ConvNeXt-Tiny}
pre-entrenado en ImageNet) con un perceptron multicapa (\textbf{MLP})
que procesa \textbf{15 morfometrias estadisticas} (Gini, $M_{20}$,
$F(G, M_{20})$, $S(G, M_{20})$, concentracion, asimetria, suavidad, Sersic,
radios de Petrosian). Entrenado con \textbf{Focal Loss}, validado con
\textbf{10-Fold Stratified CV + Test-Time Augmentation} y desplegado como
\textbf{ensemble de 10 modelos}, alcanza \textbf{F1 = 0.62 $\pm$ 0.12} y
\textbf{ROC-AUC = 0.89} sobre 1442 galaxias con severo desbalance (9.8\% clumpy).
El pipeline fue sometido a auditoria completa que detecto y corrigio
\textit{data leakage} de una version previa.
}

\column{0.35}
\block{Metricas Finales}{
\centering\large
\begin{tabular}{lr}
\toprule
\rowcolor{astroblue!20} \textbf{Metrica} & \textbf{Valor} \\
\midrule
F1-Score       & \textbf{0.62 $\pm$ 0.12} \\
ROC-AUC        & \textbf{0.89}            \\
PR-AUC         & 0.51                     \\
Precision      & 0.64                     \\
Recall         & 0.62                     \\
Accuracy       & 0.91                     \\
\bottomrule
\end{tabular}
\vspace{5mm}

\small\color{astroblue}
\textit{10-Fold CV con TTA sobre 1442 galaxias}
}
\end{columns}

% =====================================================================
% FILA 1: INTRODUCCION | METODOLOGIA | RESULTADOS
% =====================================================================
\begin{columns}

% -------- COLUMNA 1 --------
\column{0.34}
\block{1. Introduccion y Motivacion}{
\large
Las \textbf{galaxias clumpy} presentan multiples nucleos brillantes en su
morfologia, indicadores de \textbf{formacion estelar intensa} en el universo
temprano. Su clasificacion es clave para entender la evolucion galactica,
pero requiere inspeccion visual experto por experto, un proceso lento e
inconsistente entre miles de objetos catalogados.

\vspace{4mm}
\textbf{Objetivo:} Desarrollar un clasificador automatico binario
(\textit{clumpy vs no-clumpy}) que aproveche tanto las imagenes JWST como
las morfometrias estadisticas ya calculadas, superando enfoques puramente
visuales o puramente numericos.

\vspace{4mm}
\textbf{Contribuciones:}
\begin{itemize}
\item Arquitectura hibrida \textbf{imagen + tabular} (ConvNeXt + MLP) para clasificacion morfologica
\item Manejo riguroso del desbalance severo (9.8\% clase positiva) con Focal Loss + WeightedRandomSampler
\item Pipeline validado con \textbf{10-Fold CV honesto} + TTA + ensemble de 10 modelos
\item \textbf{Auditoria completa} que detecto y corrigio bugs de \textit{data leakage} en versiones previas
\item Analisis comparativo de 5 configuraciones (baseline solo-imagen, hibrido base, ResNet-50, v2 con MixUp/LS, v2.1 final)
\end{itemize}
}

\block{2. Dataset}{
\large
\begin{tabular}{ll}
\textbf{Total galaxias:}   & 1442 \\
\textbf{Clumpy:}           & \textcolor{clumpyred}{\textbf{141 (9.8\%)}} \\
\textbf{No clumpy:}        & \textcolor{noclumpygreen}{\textbf{1301 (90.2\%)}} \\
\textbf{Imagenes:}         & PNG RGB 240$\times$240 px \\
\textbf{Instrumento:}      & JWST/NIRCam \\
\textbf{Campos:}           & COS4, GS4, GN4, G4, U4 \\
\end{tabular}

\vspace{5mm}
\textbf{Etiqueta}: \texttt{visual\_clas.Clumpy} del catalogo
\texttt{11\_05\_2026\_dictionary} (clasificacion visual experta).

\vspace{5mm}
\textbf{15 morfometrias por galaxia} (statmorph):
\begin{multicols}{2}\small
\begin{itemize}\itemsep0pt
\item Gini
\item $M_{20}$
\item $F(G, M_{20})$
\item $S(G, M_{20})$
\item C (concentracion)
\item A (asimetria)
\item S (suavidad)
\item $\text{S/N}$ / pixel
\item Indice Sersic $n$
\item $r_{\text{half}}$ Sersic
\item Elipticidad
\item Elongacion
\item $r_{\text{petro,circ}}$
\item $r_{20}$, $r_{80}$
\end{itemize}
\end{multicols}
}

% -------- COLUMNA 2 --------
\column{0.33}
\block{3. Arquitectura Hibrida}{
\centering
\begin{tikzpicture}[
    node distance=6mm,
    every node/.style={font=\small\bfseries},
    box/.style={rectangle, rounded corners=5pt, draw=astroblue,
                fill=white, line width=1.5pt, minimum width=45mm,
                minimum height=12mm, align=center},
    imgbox/.style={box, fill=astrocyan!20},
    tabbox/.style={box, fill=astrogold!20},
    fusebox/.style={box, fill=noclumpygreen!20, minimum width=55mm},
    outbox/.style={box, fill=clumpyred!30, minimum width=55mm, minimum height=15mm},
    arrow/.style={-{Stealth[length=4mm]}, line width=1.8pt, color=astroblue}
]

% Entradas
\node[imgbox] (img) at (-6,4) {Imagen JWST\\ 224$\times$224$\times$3};
\node[tabbox] (tab) at (6,4) {15 morfometrias\\ (Gini, $M_{20}$, S, ...)};

% Backbones
\node[imgbox, minimum height=18mm] (convnext) at (-6,1.5) {ConvNeXt-Tiny\\ pre-entrenado ImageNet\\ (backbone congelado 5 ep)};
\node[tabbox, minimum height=18mm] (mlp) at (6,1.5) {MLP\\ Linear(15$\to$128)\\ ReLU + Dropout};

% Features
\node[box, fill=astrocyan!40] (imgfeat) at (-6,-1.5) {Features imagen\\ $\mathbb{R}^{768}$};
\node[box, fill=astrogold!40] (tabfeat) at (6,-1.5) {Features tabular\\ $\mathbb{R}^{128}$};

% Fusion
\node[fusebox] (concat) at (0,-4) {Concatenacion\\ $\mathbb{R}^{896}$};

% Classifier
\node[fusebox, fill=astrogold!30] (fc) at (0,-6.5) {Linear(896 $\to$ 1)};

% Output
\node[outbox] (out) at (0,-9.5) {\Large Clumpy / No-Clumpy\\ \small(sigmoid + umbral optimizado)};

% Arrows
\draw[arrow] (img) -- (convnext);
\draw[arrow] (tab) -- (mlp);
\draw[arrow] (convnext) -- (imgfeat);
\draw[arrow] (mlp) -- (tabfeat);
\draw[arrow] (imgfeat) -| (concat);
\draw[arrow] (tabfeat) -| (concat);
\draw[arrow] (concat) -- (fc);
\draw[arrow] (fc) -- (out);

\end{tikzpicture}

\vspace{3mm}
\raggedright\large
\textbf{Loss:} Focal Loss ($\alpha=0.9$, $\gamma=2$) para clase minoritaria.\\
\textbf{Sampler:} \texttt{WeightedRandomSampler} balancea batches en train.\\
\textbf{Optimizer:} AdamW, LR = $2\text{e-}4$ (head) $\to$ $5\text{e-}5$ (fine-tune).\\
\textbf{Scheduler:} CosineAnnealingLR con freeze de 5 epocas.
}

\block{4. Estrategia de Entrenamiento y Validacion}{
\large
\begin{itemize}
\item \textbf{10-Fold Stratified CV}: cada galaxia en exactamente 1 fold de val.
\item \textbf{Scaler por fold}: normalizacion Z-score ajustada solo en train (sin leakage).
\item \textbf{TTA (4 augmentaciones)}: original, flip-H, flip-V, flip-HV en evaluacion.
\item \textbf{Ensemble final}: promedio de los 10 modelos entrenados $\times$ 4 aug TTA = \textbf{40 predicciones por muestra}.
\item \textbf{Umbral optimo}: buscado en val (rango 0.01--0.99).
\item \textbf{Mixed precision (AMP)}: acelera en GPU T4.
\item \textbf{Sin retrain final}: se usa directamente el ensemble como modelo desplegable.
\end{itemize}

\vspace{3mm}
\small\textit{Notebook implementacion: \texttt{clumpy\_classifier\_convnext\_hybrid\_v2\_1.ipynb}
(v2.1 sin MixUp ni Label Smoothing, tras auditoria de bugs).}
}

% -------- COLUMNA 3 --------
\column{0.33}
\block{5. Resultados por Fold}{
\centering
% Reemplazar con imagen real de barras por fold
% \includegraphics[width=0.95\linewidth]{v2_1_cv_metrics.png}
\vspace{2mm}
\includegraphics[width=0.98\linewidth,keepaspectratio]{figs/cv_metrics.png}
\vspace{2mm}

\small\raggedright
\textit{F1, precision, recall y ROC-AUC por fold (10-fold CV con TTA).
Std alta en F1 (0.12) refleja la sensibilidad del split con solo 141 clumpy totales.}
}

\block{6. Curvas ROC / PR y Matriz de Confusion}{
\centering
\includegraphics[width=0.98\linewidth,keepaspectratio]{figs/cv_curves.png}

\vspace{3mm}\small\raggedright
\textit{Predicciones out-of-fold agregadas. La ROC muestra buena
discriminacion (AUC = 0.89). La PR curve refleja el desafio del desbalance:
AP = 0.51 vs baseline aleatorio = 0.098 (\textbf{5$\times$ mejor}).
La matriz de confusion usa el umbral optimo pooled.}
}

\end{columns}

% =====================================================================
% FILA 2: COMPARACION MODELOS + DISCUSION
% =====================================================================
\begin{columns}

\column{0.5}
\block{7. Comparacion de Arquitecturas Evaluadas}{
\centering\large
\begin{tabular}{lcc}
\toprule
\rowcolor{astroblue!20} \textbf{Modelo} & \textbf{F1 CV} & \textbf{ROC-AUC} \\
\midrule
ConvNeXt-Tiny (solo imagen)          & 0.58 $\pm$ 0.07 & 0.90 \\
ConvNeXt hibrido (10 feat, 5-fold)   & 0.59 $\pm$ 0.06 & 0.87 \\
\rowcolor{astrogold!20}
\textbf{ConvNeXt hibrido v2.1} \textbf{$\star$}       & \textbf{0.62 $\pm$ 0.12} & \textbf{0.89} \\
\quad (15 feat + TTA + 10-fold)      &                 &      \\
ResNet-50 hibrido                    & 0.50            & 0.86 \\
v2 (con MixUp + LS)                  & 0.58 $\pm$ 0.08 & 0.89 \\
\bottomrule
\end{tabular}

\vspace{5mm}
\raggedright\large\textbf{Auditoria y correccion de bugs:}
Durante el desarrollo se identifico un \textbf{bug de \textit{data leakage}} en
la version base del hibrido (unpacking incorrecto de \texttt{make\_loaders}
que hacia evaluar sobre \texttt{train\_loader} con oversampling). Esto inflaba
artificialmente el F1 a 0.76. Tras el fix, todos los modelos fueron
re-validados con metricas honestas.

\vspace{3mm}
\textbf{v2.1} es la version final: elimina el retrain sobre samples completo
(evita segundo leakage), usa \textbf{15 features tabulares} (agrega
$F(G, M_{20})$, $S(G, M_{20})$, $r_{\text{petro,circ}}$, $r_{20}$, $r_{80}$),
y aplica \textbf{TTA en cada fold} para reducir varianza en evaluacion.
Ademas se auditaron y corrigieron bugs menores: \texttt{best\_state}
inicializado, rango de umbral ampliado (0.01--0.99), y guard clause para
folds con clase unica.
}

\column{0.5}
\block{8. Discusion}{
\large
\textbf{Fortalezas:}
\begin{itemize}
\item Discriminacion solida (ROC-AUC = 0.89): el modelo ordena correctamente
las galaxias por probabilidad de ser clumpy.
\item Pipeline robusto sin \textit{data leakage}: scaler por fold, val disjunta,
umbral optimizado en val (no en test).
\item El componente hibrido \textbf{imagen + tabular} aporta marginalmente
sobre solo-imagen (F1 = 0.62 vs 0.58), sugiriendo que ConvNeXt ya captura
buena parte de las morfometrias implicitamente.
\end{itemize}

\vspace{4mm}
\textbf{Limitaciones:}
\begin{itemize}
\item \textbf{Techo por tamano del dataset}: multiples arquitecturas
(ConvNeXt, ResNet-50, EfficientNet) convergen a F1 = 0.50--0.67 con solo
141 ejemplos positivos. La limitacion es de datos, no de modelo.
\item \textbf{Precision moderada} (0.64): $\sim$1 de cada 2 predicciones
``clumpy'' es correcta; insuficiente para pipelines totalmente automatizados.
\item \textbf{Variabilidad de recall} entre folds ($\sigma = 0.16$):
sensible a que galaxias clumpy caen en val.
\end{itemize}

\vspace{4mm}
\textbf{Comparacion con literatura:}
Nuestro F1 = 0.62 esta en el rango tipico reportado para clasificacion binaria
de morfologias raras con datasets similares (F1 = 0.50--0.75). Ensemble
multi-arquitectura + TTA en state-of-art alcanza F1 = 0.65--0.75.
}

\end{columns}

% =====================================================================
% FILA FINAL: CONCLUSIONES + FUTURO + REFERENCIAS + QR
% =====================================================================
\begin{columns}

\column{0.5}
\block{9. Conclusiones y Trabajo Futuro}{
\large
\begin{itemize}
\item El clasificador hibrido \textbf{ConvNeXt-Tiny + MLP} logra
\textbf{F1 = 0.62 $\pm$ 0.12} con \textbf{ROC-AUC = 0.89} en validacion
cruzada honesta, siendo competitivo con la literatura para tareas
similares con desbalance severo.
\item El modelo aprende \textbf{morfologia diferencial} y no colapsa a
la clase mayoritaria (recall = 0.62 sobre clase minoritaria = 6$\times$ mejor que baseline).
\item Las morfometrias estadisticas aportan poca informacion adicional
al backbone convolucional moderno.
\item La \textbf{auditoria de bugs} fue clave: el pipeline inicial reportaba
F1 = 0.76 (inflado por leakage). El F1 real es 0.62 tras las correcciones.
\end{itemize}

\vspace{4mm}
\textbf{Direcciones futuras:}
\begin{itemize}
\item \textbf{Pseudo-labeling}: entrenar sobre galaxias no etiquetadas usando
predicciones de alta confianza.
\item \textbf{Pre-entrenamiento self-supervised} sobre imagenes JWST no
etiquetadas (SimCLR, DINO, MAE) para mejor inicializacion.
\item \textbf{Ensemble multi-arquitectura}: combinar ConvNeXt +
EfficientNet + Vision Transformer.
\item \textbf{Anotacion adicional}: expandir el conjunto etiquetado con
mas expertos y validacion espectroscopica.
\item \textbf{Segmentacion de clumps}: pasar de clasificacion binaria a
localizacion de clumps individuales dentro de cada galaxia.
\end{itemize}
}

\column{0.5}
\block{Referencias}{
\small
\begin{enumerate}\itemsep2pt
\item Liu, Z. et al. (2022). \textit{A ConvNet for the 2020s (ConvNeXt)}.
CVPR 2022.
\item Lin, T-Y. et al. (2017). \textit{Focal Loss for Dense Object Detection}.
ICCV 2017.
\item Rodriguez-Gomez, V. et al. (2019). \textit{statmorph: Python code for
calculating non-parametric morphological diagnostics of galaxy images}.
MNRAS 483:4140.
\item Lotz, J. et al. (2008). \textit{Galaxy morphology in the Gini-M20 plane}.
MNRAS 391:1137.
\item Deng, J. et al. (2009). \textit{ImageNet: A Large-Scale Hierarchical
Image Database}. CVPR 2009.
\item Gardner, J. et al. (2023). \textit{The James Webb Space Telescope
Mission}. PASP 135:1048.
\item Wuyts, S. et al. (2012). \textit{Smooth(er) Stellar Mass Maps in CANDELS:
Constraints on the Longevity of Clumps in High-Redshift Star-Forming Galaxies}.
ApJ 753:114.
\end{enumerate}
}

\block{Codigo, Datos y Reproducibilidad}{
\centering
\begin{minipage}[c]{0.35\linewidth}
\centering
\qrcode[height=6cm]{https://github.com/anomalyco/Proyecto-Semestral-Deep-Learning}\\[3mm]
\small\textbf{Repositorio GitHub}
\end{minipage}
\hspace{5mm}
\begin{minipage}[c]{0.55\linewidth}
\large
\textbf{Stack:} PyTorch 2.6 $|$ timm 1.0 $|$ scikit-learn 1.3 $|$ Google Colab (T4).\\[3mm]
\textbf{Notebook final:} \\
\texttt{clumpy\_classifier\_convnext\_hybrid\_v2\_1.ipynb}\\[3mm]
\textbf{Modelo desplegable:} \texttt{ensemble\_clumpy\_v2\_1.pth}\\
(10 ConvNeXt-Tiny + scalers + umbral, 15 features tabulares)\\[3mm]
\textbf{Dataset:} 1442 galaxias JWST etiquetadas manualmente\\
(catalogo \texttt{11\_05\_2026\_dictionary\_with\_all\_statmorph}).

\vspace{3mm}
Escanea el QR o visita el repositorio para reproducir todos los experimentos.

\vspace{3mm}
\small\textit{Agradecimientos: al equipo de anotacion del catalogo JWST y
al programa Google Colab por acceso a GPU T4.}
\end{minipage}
}

\end{columns}

\end{document}
